\documentclass[11pt,a4paper]{article}

\usepackage[margin=2.5cm]{geometry}
\usepackage[hidelinks]{hyperref}
\usepackage{listings}
\usepackage{xcolor}
\usepackage{booktabs}
\usepackage{array}
\usepackage{enumitem}
\usepackage{parskip}
\usepackage{titlesec}
\usepackage{caption}
\usepackage{subcaption}
\usepackage{graphicx}
\usepackage{multirow}

% ── Listing styles ────────────────────────────────────────────────────────────
\definecolor{codebg}{rgb}{0.97,0.97,0.97}
\definecolor{kw}{rgb}{0.13,0.29,0.53}
\definecolor{cmt}{rgb}{0.35,0.50,0.35}
\definecolor{str}{rgb}{0.65,0.12,0.12}

\lstdefinestyle{sparql}{
  backgroundcolor=\color{codebg},
  basicstyle=\ttfamily\small,
  keywordstyle=\color{kw}\bfseries,
  commentstyle=\color{cmt}\itshape,
  stringstyle=\color{str},
  breaklines=true,
  frame=single,
  rulecolor=\color{gray!40},
  numbers=none,
  morekeywords={SELECT,WHERE,PREFIX,DISTINCT,FILTER,OPTIONAL,UNION,FROM,ASK,CONSTRUCT},
}

\lstdefinestyle{turtle}{
  backgroundcolor=\color{codebg},
  basicstyle=\ttfamily\small,
  keywordstyle=\color{kw}\bfseries,
  commentstyle=\color{cmt}\itshape,
  stringstyle=\color{str},
  breaklines=true,
  frame=single,
  rulecolor=\color{gray!40},
  numbers=none,
  morekeywords={@prefix,a},
}

\lstdefinestyle{python}{
  backgroundcolor=\color{codebg},
  basicstyle=\ttfamily\small,
  keywordstyle=\color{kw}\bfseries,
  commentstyle=\color{cmt}\itshape,
  stringstyle=\color{str},
  breaklines=true,
  frame=single,
  rulecolor=\color{gray!40},
  numbers=left,
  numberstyle=\tiny\color{gray},
  language=Python,
}

\lstdefinestyle{cypher}{
  backgroundcolor=\color{codebg},
  basicstyle=\ttfamily\small,
  keywordstyle=\color{kw}\bfseries,
  commentstyle=\color{cmt}\itshape,
  stringstyle=\color{str},
  breaklines=true,
  frame=single,
  rulecolor=\color{gray!40},
  numbers=none,
  morekeywords={MATCH,RETURN,WHERE,WITH,CREATE,MERGE,SET,DELETE,OPTIONAL,UNION,AS,CALL,IN,IS,NOT,NULL,AND,OR},
}

\begin{document}

% -- Title
\pagestyle{empty}
\input{portada}

% -- Contents
 
\tableofcontents
\newpage

% -- Document
\pagestyle{plain}
\setcounter{page}{1}


\section{Section A — Modeling}


\subsection{Overview}

This section describes the RDFS knowledge graph designed to model relationships
between drugs and diseases in a biomedical domain.  The graph is organized in
three canonical layers:

\begin{enumerate}[noitemsep]
  \item \textbf{Metamodel layer} — the RDF/RDFS vocabulary constructs used to
        define classes and properties.
  \item \textbf{Model layer} — the domain taxonomy (class hierarchy and
        property hierarchy).
  \item \textbf{Facts layer} — the instances (individual drugs and diseases) and
        the explicit relationships between them.
\end{enumerate}

The two root entities are \texttt{bio:Drug} and \texttt{bio:Disease}.  Drugs are
specialized into \emph{AntiInflammatory}, \emph{Antibiotic}, \emph{Analgesic},
\emph{Antiviral}, and \emph{Steroid}.  Diseases are specialized into
\emph{InflammatoryDisease}, \emph{InfectiousDisease}, \emph{ChronicDisease},
and \emph{NeurologicalDisease}.  A single top-level property \texttt{bio:affects}
is refined into three subproperties: \texttt{treats}, \texttt{relieves}, and
\texttt{worsens}.

The namespace used throughout is:
\begin{center}
  \texttt{http://www.example.org/biomedical\#} \quad (prefix \texttt{bio:})
\end{center}

Here is a picture of the metamodel and model layers, as well as one example of an instance:

\begin{figure}[h!]
    \centering
    \includegraphics[width=0.65\textwidth]{imagenes/modeling.png}
    \caption{Metamodel, model and instance layers of the graph}
    \label{fig:graph}
\end{figure}


% ── A.2 ──────────────────────────────────────────────────────────────────────
\subsection{Metamodel Layer}

The metamodel layer uses only standard RDF/RDFS vocabulary.  It expresses
\emph{how classes and properties are constructed}, independently of the domain
content.  The key constructs employed are listed in
Table~\ref{tab:metamodel}.

\begin{table}[h!]
\centering
\caption{RDF/RDFS constructs used in the metamodel layer.}
\label{tab:metamodel}
\begin{tabular}{@{}lp{9cm}@{}}
\toprule
\textbf{Construct} & \textbf{Role in this graph} \\
\midrule
\texttt{rdfs:Class}           & Declares \texttt{bio:Drug}, \texttt{bio:Disease}, and all subclasses as RDF classes. \\
\texttt{rdfs:subClassOf}      & Establishes the class hierarchy (e.g., \texttt{AntiInflammatory rdfs:subClassOf Drug}). \\
\texttt{rdf:Property}         & Declares \texttt{bio:affects}, \texttt{bio:treats}, \texttt{bio:relieves}, \texttt{bio:worsens} as properties. \\
\texttt{rdfs:subPropertyOf}   & Establishes the property hierarchy (e.g., \texttt{treats rdfs:subPropertyOf affects}). \\
\texttt{rdfs:domain}          & Constrains the subject of \texttt{affects} to be a \texttt{Drug}. \\
\texttt{rdfs:range}           & Constrains the object of \texttt{affects} to be a \texttt{Disease}. \\
\texttt{rdf:type}             & Assigns each instance to exactly one leaf class (e.g., \texttt{Ibuprofen rdf:type AntiInflammatory}). \\
\bottomrule
\end{tabular}
\end{table}

\noindent Note that \texttt{rdfs:domain} and \texttt{rdfs:range} are declared
only on \texttt{bio:affects}; the subproperties inherit them through the
property hierarchy.  Likewise, \texttt{rdf:type} arrows from model-layer
elements to metamodel-layer elements (e.g.\ \texttt{bio:Drug rdf:type
rdfs:Class}) are not asserted explicitly — they are derived from RDFS
built-in axioms.  Both decisions are justified in
Section~\ref{subsec:design}.

The metamodel layer is illustrated in the top section of Figure~\ref{fig:graph}.

% ── A.3 ──────────────────────────────────────────────────────────────────────
\subsection{Model Layer}

The model layer defines the \emph{domain-specific} taxonomy: which classes
exist and how they are related, and which properties constitute the vocabulary.

\subsubsection*{Class Hierarchy}

\begin{verbatim}
bio:Drug
    bio:AntiInflammatory
    bio:Antibiotic

bio:Disease
    bio:InflammatoryDisease
    bio:InfectiousDisease
\end{verbatim}

Each subclass is declared with \texttt{rdfs:subClassOf} pointing to its
immediate parent.  The implementation extends this base schema with additional
subclasses to reach 100 instances; see Section~\ref{sec:implementation} for
details.

\subsubsection*{Property Hierarchy}

\begin{verbatim}
bio:affects          (domain: bio:Drug  |  range: bio:Disease)
    bio:treats       -- subPropertyOf bio:affects
    bio:relieves     -- subPropertyOf bio:affects
    bio:worsens      -- subPropertyOf bio:affects
\end{verbatim}



% ── A.4 ──────────────────────────────────────────────────────────────────────
\subsection{Facts Layer}

\subsubsection*{Instances}

The graph contains exactly \textbf{100 instances}: 50 drugs spread across the
five drug subclasses, and 50 diseases spread across the four disease subclasses.
Table~\ref{tab:instances} summarizes the counts.

\begin{table}[h!]
\centering
\caption{Instance distribution across leaf classes.}
\label{tab:instances}
\begin{tabular}{@{}llr@{}}
\toprule
\textbf{Root class} & \textbf{Subclass} & \textbf{Count} \\
\midrule
\multirow{5}{*}{Drug}    & AntiInflammatory & 10 \\
                         & Antibiotic       & 10 \\
                         & Analgesic        & 10 \\
                         & Antiviral        & 10 \\
                         & Steroid          & 10 \\
\midrule
\multirow{4}{*}{Disease} & InflammatoryDisease  & 13 \\
                         & InfectiousDisease    & 13 \\
                         & ChronicDisease       & 12 \\
                         & NeurologicalDisease  & 12 \\
\midrule
\multicolumn{2}{l}{\textbf{Total}} & \textbf{100} \\
\bottomrule
\end{tabular}
\end{table}

\subsubsection*{Explicit Relationship Triples}

Each instance is typed to exactly one leaf class using \texttt{rdf:type}.
Relationship triples use only the three specific subproperties (\texttt{treats},
\texttt{relieves}, \texttt{worsens}) — the \texttt{affects} property is
\emph{never asserted explicitly} between instances.  GraphDB's RDFS reasoner
will derive all \texttt{affects} triples automatically via the
\texttt{subPropertyOf} rule.  A sample of the explicit facts is shown below:

\begin{lstlisting}[style=turtle, caption={Sample Turtle triples from the facts layer.}]
bio:Ibuprofen       rdf:type          bio:AntiInflammatory .
bio:Ibuprofen       bio:relieves      bio:Arthritis .
bio:Ibuprofen       bio:worsens       bio:ChronicKidneyDisease 
bio:Amoxicillin     rdf:type          bio:Antibiotic .
bio:Amoxicillin     bio:treats        bio:Pneumonia .
\end{lstlisting}

In total the graph contains 76 explicit relationship triples (29 \texttt{treats},
33 \texttt{relieves}, 14 \texttt{worsens}) and 205 triples overall including the
metamodel and model declarations.

% ── A.5 ──────────────────────────────────────────────────────────────────────
\subsection{Design Justifications}
\label{subsec:design}

\paragraph{Why a property hierarchy for \texttt{affects}?}
A single abstract property \texttt{affects} unifies the three specific
relationships.  This lets users write a single
SPARQL query for "all drug–disease interactions" without knowing 
whether the relationship is therapeutic, symptomatic, or contraindicated.
Without this hierarchy, a query would need a \texttt{UNION} of three triple
patterns; with it, querying \texttt{bio:affects} suffices.

\paragraph{Why minimise explicit triples?}
To rely on RDFS inference rather
than asserting every derivable triple. By never asserting
\texttt{bio:Ibuprofen rdf:type bio:Drug} explicitly; GraphDB derives it from
\texttt{Ibuprofen rdf:type AntiInflammatory} and
\texttt{AntiInflammatory rdfs:subClassOf Drug}.
This approach reduces graph size, eliminates the risk of inconsistent redundancies.

\paragraph{Why \texttt{rdfs:domain} and \texttt{rdfs:range} only on \texttt{affects}?}
Declaring constraints at the most general level and relying on inheritance
prevents repeating them on each subproperty.  RDFS does not mandate that
subproperties inherit domain/range, but GraphDB's RDFS (Optimized) ruleset
does apply the standard RDFS-entailment rules (rdfs4, rdfs7, rdfs9) so this
works as expected.



% ── A.6 ──────────────────────────────────────────────────────────────────────
\subsection{RDFS Knowledge Graph vs.\ Neo4j Property Graph}

\noindent \textbf{Key gap}: the capability that RDFS provides and Neo4j does not is
\emph{ontological inference} — the ability to derive new facts from structural
constraints (subclass, subproperty, domain, range) without rewriting queries
or duplicating data.  A Neo4j property graph can store the same raw facts, but
it cannot answer ``find all drugs that affect any disease'' without a
programmer explicitly enumerating every relationship type in the query.
Explicit \texttt{rdfs:subClassOf} chain; querying a parent class returns all
  subclass instances automatically.
Node labels are flat; no inheritance.  A query for \texttt{Drug} returns only
  nodes with that exact label — not \texttt{Antibiotic} nodes unless explicitly
  listed.
Facts are stated once at the most specific level; all generalizations are inferred. 
Every fact that should be queryable must be explicitly present.  A new drug class requires updating all queries that should match it.




% ─────────────────────────────────────────────────────────────────────────────
\newpage
\section{Section B — Implementation}
\label{sec:implementation}
% ─────────────────────────────────────────────────────────────────────────────

The base schema defined in Section~A contains two drug subclasses
(\texttt{AntiInflammatory}, \texttt{Antibiotic}) and two disease subclasses
(\texttt{InflammatoryDisease}, \texttt{InfectiousDisease}).  To reach the
required 100 instances, the implementation extends the schema with three
additional drug subclasses — \texttt{Analgesic}, \texttt{Antiviral},
\texttt{Steroid} — and two additional disease subclasses —
\texttt{ChronicDisease}, \texttt{NeurologicalDisease}.  All five new subclasses
are declared with \texttt{rdfs:subClassOf} pointing to their respective root
class, following the same pattern as the original schema.

The knowledge graph has five logical blocks:

\begin{enumerate}[noitemsep]
  \item \textbf{Metamodel declarations} — \texttt{rdfs:Class},
        \texttt{rdfs:subClassOf}, \texttt{rdf:Property},
        \texttt{rdfs:subPropertyOf}, \texttt{rdfs:domain}, \texttt{rdfs:range}.
  \item \textbf{Instance typing} — one \texttt{rdf:type} triple per instance
        pointing to the most specific (leaf) class.
  \item \textbf{Relationship triples} — \texttt{treats}, \texttt{relieves},
        \texttt{worsens} facts derived from clinical knowledge.
  \item \textbf{Serialisation} — \texttt{g.serialize(format="turtle")} writes
        \texttt{biomedical\_kg.ttl}.
  \item \textbf{GraphDB upload} — waits for the server, creates the repository
        using a Turtle configuration file posted as a multipart form to
        \texttt{/rest/repositories}, then uploads the triples via
        \texttt{POST /repositories/biomedical/statements}.
\end{enumerate}

We created the script by  the \textbf{RDFS (Optimized)} ruleset
(\texttt{rdfsplus-optimized}), which is GraphDB's full RDFS entailment regime
that materialises inferred triples at load time.

We added 100 instances that follow triple counts as it is seen in Table~\ref{tab:counts}
\begin{table}[h!]
\centering
\caption{Triple counts in the generated graph.}
\label{tab:counts}
\begin{tabular}{@{}lr@{}}
\toprule
\textbf{Category} & \textbf{Count} \\
\midrule
Drug instances (\texttt{rdf:type} leaf class)    & 50 \\
Disease instances (\texttt{rdf:type} leaf class) & 50 \\
\texttt{treats} relationships                    & 29 \\
\texttt{relieves} relationships                  & 33 \\
\texttt{worsens} relationships                   & 14 \\
Metamodel and model declarations                 & 29 \\
\midrule
\textbf{Total triples in .ttl file}              & \textbf{205} \\
\bottomrule
\end{tabular}
\end{table}

% ─────────────────────────────────────────────────────────────────────────────
\newpage
\section{Section C — Querying}
% ─────────────────────────────────────────────────────────────────────────────

\subsection{Assumptions}

\begin{itemize}[noitemsep]
  \item GraphDB is running with the \textbf{RDFS (Optimized)} ruleset
        (\texttt{rdfsplus-optimized}), which materialises inferred triples at
        load time so that SPARQL queries over the default graph see both
        explicit and inferred triples.
  \item The equivalent Neo4j graph stores the same 100 nodes and 76
        relationships, but uses flat labels (\texttt{:Drug},
        \texttt{:AntiInflammatory}, \ldots) and typed relationships
        (\texttt{:TREATS}, \texttt{:RELIEVES}, \texttt{:WORSENS}) with no
        schema hierarchy.  Every node that should match a ``drug'' query carries
        both the specific label (e.g.\ \texttt{:Antibiotic}) and the general
        label (\texttt{:Drug}) — these must be added explicitly by the data
        loader; they are not derived.
  \item SPARQL queries use the namespace prefix
        \texttt{bio: <http://www.example.org/biomedical\#>}.
  \item \texttt{SELECT DISTINCT} is used throughout to suppress any duplicate
        bindings that could arise from multiple inferred paths to the same triple.
\end{itemize}

% ── Query 1 ───────────────────────────────────────────────────────────────────
\subsection{Query 1 — List All Drugs}

\begin{lstlisting}[style=sparql, caption={SPARQL: list all drugs.}]
PREFIX bio: <http://www.example.org/biomedical#>
PREFIX rdf: <http://www.w3.org/1999/02/22-rdf-syntax-ns#>

SELECT DISTINCT ?drug
WHERE {
    ?drug rdf:type bio:Drug .
}
ORDER BY ?drug
\end{lstlisting}

\begin{lstlisting}[style=cypher, caption={Cypher equivalent: list all drugs.}]
MATCH (d:Drug)
RETURN d.name AS drug
ORDER BY drug
\end{lstlisting}

\subsubsection*{Comparison with Neo4j}

The surface syntax is very similar, but the \emph{semantics} differ
fundamentally.  In the SPARQL version, \texttt{bio:Drug} is a class in an RDFS
hierarchy; querying it returns all instances of any subclass without listing
them.  In the Cypher version, \texttt{:Drug} is a plain label — it only
matches nodes that carry exactly that label. The Cypher query returns all 50 drug nodes because each was manually labeled :Drug. If a new drug subclass is added later, the query will only work if the :Drug label is also applied — but Neo4j has no schema mechanism to enforce this.  In SPARQL, adding a new subclass
(\texttt{bio:Antifungal rdfs:subClassOf bio:Drug}) is sufficient; the query
needs no modification.

\subsubsection*{Inferences applied by GraphDB}

The query pattern \texttt{?drug rdf:type bio:Drug} triggers RDFS rule \textbf{rdfs9},
which states: if \texttt{?u rdf:type ?C} and \texttt{?C rdfs:subClassOf ?D}, then
infer \texttt{?u rdf:type ?D}. None of the 50 drug instances are explicitly typed
as \texttt{bio:Drug} in the uploaded Turtle file; they are typed only as their leaf
class. GraphDB materialises the following inferred triples (one per instance, five
subclass chains):

\begin{lstlisting}[style=turtle]
bio:Ibuprofen        rdf:type  bio:Drug .  # via AntiInflammatory subClassOf Drug
bio:Amoxicillin      rdf:type  bio:Drug .  # via Antibiotic       subClassOf Drug
bio:Paracetamol      rdf:type  bio:Drug .  # via Analgesic        subClassOf Drug
# ... 45 more (one per remaining instance)
\end{lstlisting}


% ── Query 2 ───────────────────────────────────────────────────────────────────
\subsection{Query 2 — List All Diseases}

\begin{lstlisting}[style=sparql, caption={SPARQL: list all diseases.}]
PREFIX bio: <http://www.example.org/biomedical#>
PREFIX rdf: <http://www.w3.org/1999/02/22-rdf-syntax-ns#>

SELECT DISTINCT ?disease
WHERE {
    ?disease rdf:type bio:Disease .
}
ORDER BY ?disease
\end{lstlisting}

\begin{lstlisting}[style=cypher, caption={Cypher equivalent: list all diseases.}]
MATCH (d:Disease)
RETURN d.name AS disease
ORDER BY disease
\end{lstlisting}

\subsubsection*{Comparison with Neo4j}

Identical structural argument as Query 1.  The Cypher query works only because
every disease node was given the \texttt{:Disease} label in addition to its
specific label (\texttt{:InflammatoryDisease}, etc.) at load time.  The SPARQL
query works because RDFS subclass inference makes it unnecessary to assert the
general type explicitly.  The two queries produce the same result set, but
one relies on inference while the other relies on redundant labelling.

\subsubsection*{Inferences applied by GraphDB}

The same rdfs9 rule applies across four subclass chains:

\begin{lstlisting}[style=turtle]
bio:Arthritis          rdf:type  bio:Disease .  # via InflammatoryDisease subClassOf Disease
bio:Pneumonia          rdf:type  bio:Disease .  # via InfectiousDisease   subClassOf Disease
bio:Diabetes           rdf:type  bio:Disease .  # via ChronicDisease      subClassOf Disease
# ... 46 more
\end{lstlisting}


% ── Query 3 ───────────────────────────────────────────────────────────────────
\subsection{Query 3 — Drug–Disease Pairs Where Drug Affects Disease}

\begin{lstlisting}[style=sparql, caption={SPARQL: all drug--disease interaction pairs.}]
PREFIX bio: <http://www.example.org/biomedical#>

SELECT DISTINCT ?drug ?disease
WHERE {
    ?drug bio:affects ?disease .
}
ORDER BY ?drug ?disease
\end{lstlisting}

\begin{lstlisting}[style=cypher, caption={Cypher equivalent: all drug--disease interactions.}]
MATCH (d)-[r:TREATS|RELIEVES|WORSENS]->(dis)
RETURN d.name AS drug, dis.name AS disease
ORDER BY drug, disease
\end{lstlisting}

\subsubsection*{Comparison with Neo4j}

This is the query where the two models diverge most sharply.
In SPARQL, a single triple pattern on \texttt{bio:affects} retrieves all 76
pairs — \texttt{bio:affects} is never explicitly used in the data, but RDFS
infers it from the three subproperties.  In Cypher, there is no concept of a
property superset; the query must enumerate all three relationship types
using the \texttt{|} union syntax.  If a fourth relationship type (e.g.\
\texttt{:PREVENTS}) were added, the Cypher query would silently omit it unless
the developer remembered to update the pattern.  The SPARQL query would include
it automatically, provided only that the new property is declared
\texttt{rdfs:subPropertyOf bio:affects}.

A further difference is expressiveness: Neo4j has no way to write a query
equivalent to ``find all drug–disease pairs connected by \emph{any} relationship
that is a subtype of \texttt{affects}'' without explicitly naming the subtypes.
RDFS makes this possible through a single abstract property.

\subsubsection*{Inferences applied by GraphDB}

Two rules cooperate here.

\textbf{Rule rdfs7} (subproperty propagation): Applied to every explicit relationship triple, this produces 76 inferred
\texttt{bio:affects} triples — one for each asserted
\texttt{treats}/\texttt{relieves}/\texttt{worsens} fact.  Examples:

\begin{lstlisting}[style=turtle]
# Explicit:
bio:Ibuprofen    bio:relieves  bio:Arthritis .
bio:Amoxicillin  bio:treats    bio:Pneumonia .
bio:Ibuprofen    bio:worsens   bio:ChronicKidneyDisease .

# Inferred via rdfs7:
bio:Ibuprofen    bio:affects   bio:Arthritis .
bio:Amoxicillin  bio:affects   bio:Pneumonia .
bio:Ibuprofen    bio:affects   bio:ChronicKidneyDisease .
\end{lstlisting}

\textbf{Rules rdfs2 and rdfs3} (domain and range typing): since
\texttt{bio:affects} has \texttt{rdfs:domain bio:Drug} and
\texttt{rdfs:range bio:Disease}, GraphDB additionally materialises:

\begin{lstlisting}[style=turtle]
# For every subject of an affects triple (via rdfs2):
bio:Ibuprofen    rdf:type  bio:Drug    .
bio:Amoxicillin  rdf:type  bio:Drug    .

# For every object of an affects triple (via rdfs3):
bio:Arthritis    rdf:type  bio:Disease .
bio:Pneumonia    rdf:type  bio:Disease .
\end{lstlisting}

Note that these domain/range inferences \emph{overlap} with the rdfs9 subclass
inferences from Queries 1 and 2.  Both paths produce the same
\texttt{rdf:type bio:Drug} and \texttt{rdf:type bio:Disease} triples, which
GraphDB deduplicates.  This redundancy is harmless and actually desirable: it
means the typing is robust even if one inference path were to fail.

\newpage

% ── Summary ───────────────────────────────────────────────────────────────────
\subsection{Summary of Inferred Triple Patterns}

\begin{table}[h!]
\centering
\caption{All inferred triple patterns, the RDFS rules that produce them, and a
         concrete example for each.}
\renewcommand{\arraystretch}{1.35}
\begin{tabular}{@{}p{4.6cm}p{1.8cm}p{7cm}@{}}
\toprule
\textbf{Inferred pattern} & \textbf{Rule} & \textbf{Concrete example} \\
\midrule
\texttt{?x rdf:type bio:Drug} &
rdfs9 &
\texttt{Ibuprofen rdf:type Drug} derived from
\texttt{Ibuprofen rdf:type AntiInflammatory} and
\texttt{AntiInflammatory subClassOf Drug}. \\

\texttt{?x rdf:type bio:Disease} &
rdfs9 &
\texttt{Arthritis rdf:type Disease} derived from
\texttt{Arthritis rdf:type InflammatoryDisease} and
\texttt{InflammatoryDisease subClassOf Disease}. \\

\texttt{?d bio:affects ?dis} &
rdfs7 &
\texttt{Ibuprofen affects Arthritis} derived from
\texttt{Ibuprofen relieves Arthritis} and
\texttt{relieves subPropertyOf affects}. \\

\texttt{?d rdf:type bio:Drug} (from domain) &
rdfs2 &
\texttt{Amoxicillin rdf:type Drug} derived from
\texttt{Amoxicillin treats Pneumonia} and
\texttt{affects rdfs:domain Drug} (cascades via rdfs7). \\

\texttt{?dis rdf:type bio:Disease} (from range) &
rdfs3 &
\texttt{Pneumonia rdf:type Disease} derived from
\texttt{Amoxicillin treats Pneumonia} and
\texttt{affects rdfs:range Disease}. \\
\bottomrule
\end{tabular}
\end{table}

\noindent All five patterns behave as expected.  The RDFS (Optimized) ruleset
in GraphDB implements the complete RDFS entailment regime.  Every inferred triple
is materialised into the repository before query evaluation begins, so query
performance is not affected by the number of inference steps required.

\end{document}
